% ****** Start of file apssamp.tex ******
%
%   This file is part of the APS files in the REVTeX 4.2 distribution.
%   Version 4.2a of REVTeX, December 2014
%
%   Copyright (c) 2014 The American Physical Society.
%
%   See the REVTeX 4 README file for restrictions and more information.
%
% TeX'ing this file requires that you have AMS-LaTeX 2.0 installed
% as well as the rest of the prerequisites for REVTeX 4.2
%
% See the REVTeX 4 README file
% It also requires running BibTeX. The commands are as follows:
%
%  1)  latex apssamp.tex
%  2)  bibtex apssamp
%  3)  latex apssamp.tex
%  4)  latex apssamp.tex
%
\documentclass[%
 reprint,
%superscriptaddress,
%groupedaddress,
%unsortedaddress,
%runinaddress,
%frontmatterverbose, 
%preprint,
%preprintnumbers,
%nofootinbib,
%nobibnotes,
%bibnotes,
 amsmath,amssymb,
 aps,
%pra,
%prb,
%rmp,
%prstab,
%prstper,
%floatfix,
]{revtex4-2}

\usepackage{graphicx}% Include figure files
\graphicspath{{figures/}} % Location of the graphics files
\usepackage{dcolumn}% Align table columns on decimal point
\usepackage{bm}% bold math
%\usepackage{hyperref}% add hypertext capabilities
%\usepackage[mathlines]{lineno}% Enable numbering of text and display math
%\linenumbers\relax % Commence numbering lines

%\usepackage[showframe,%Uncomment any one of the following lines to test 
%%scale=0.7, marginratio={1:1, 2:3}, ignoreall,% default settings
%%text={7in,10in},centering,
%%margin=1.5in,
%%total={6.5in,8.75in}, top=1.2in, left=0.9in, includefoot,
%%height=10in,a5paper,hmargin={3cm,0.8in},
%]{geometry}

\newcommand{\Bzero}{\mathbf{B}_0}
\newcommand{\Uzero}{\mathbf{U}_0}
\newcommand{\vecpot}{\mathbf{A}}
\newcommand{\linop}{\mathcal{L}}
\renewcommand{\Re}{\mathrm{Re}}

\begin{document}

\preprint{APS/123-QED}

\title{Non-modal growth and optimal perturbations in magnetohydrodynamic shear flows}

\author{Adrian Fraser}
\email{adrian_fraser@colorado.edu}
\affiliation{%
Department of Astrophysical and Planetary Sciences, University of Colorado, Boulder, CO
}%

\author{Alexis K. Kaminski}
\affiliation{
Department of Mechanical Engineering, University of California, Berkeley, Berkeley, CA
}%
\author{Jeffrey S. Oishi}
\affiliation{
Department of Physics and Astronomy, Bates College, Lewiston, ME
}%

\date{\today}% It is always \today, today,
             %  but any date may be explicitly specified

\begin{abstract}
Shear layers, MHD, optimal perturbations.
\end{abstract}

%\keywords{Suggested keywords}%Use showkeys class option if keyword
                              %display desired
\maketitle

%\tableofcontents

\section{Introduction}
    \begin{itemize}
        \item Shear flows are important in astro, space, and fusion plasmas
        \item In these cases, the system is non-normal, leading to nonmodal effects that have been thoroughly explored in the fluids community and been shown to matter but haven't received as much attention in plasmas
        \item Mention that both stratification and magnetic fields stabilize KH, to set the stage for contrasting these two systems
        \item Here or in the body, worth mentioning lift-up and Orr -- because one cool result is how lift-up is unimpeded here
        \item I think we want to find a way to connect to Squire's theorem, and how it's pretty widely invoked in plasmas to motivate 2D studies even though it's been known not to hold in MHD since the 1950's -- we want to set the stage for demonstrating that nonmodal growth is a serious thing that you're leaving out of your calculations when restricting to 2D
    \end{itemize}
[The following two paragraphs were written by Jeff, I think, at least a year ago] Shear layers are of paramount importance in numerous magnetohydrodynamic flows, in both astrophysical (cite) and fusion (cite) settings. The stability and transition to turbulence of such configurations has attracted considerable attention. In such flows, the linear operator $\linop$ is non-normal, leading to the possibility that a given perturbation can undergo significant amplification even when the flow is linearly stable. This behavior is well known in the hydrodynamics community, where it plays a pivotal role in the transition to turbulence in both canonical flows (e.g. pipe flow) and in more complex applications (e.g. oceanic mixing). While there  have been some prior studies incorporating non-modal growth in MHD flows (cite Squire, others), 

Here, we consider the linear dynamics of a simplified shear layer with a stabilizing magnetic field $\Bzero$. This system can be unstable to Kelvin-Helmholtz instabilities at weak field strength. However, even when the field is strong enough stabilize the flow, the non-modal character of the linear operator leads to significant transient growth. We compute the optimal linear perturbations that maximize that growth.

For nonzero streamwise wavenumbers that are comparable to those typically used in simulations of stratified and/or magnetized shear flows, we find results that are broadly consistent with the stratified case [cite Alexis]. When the wavenumber becomes much smaller, however, we find [etc.]

%%%%%%%%%%%%%%%%%%%%%%%%%%%%%%%%%%%%
%%%%%%%%%%%%%%%%%%%%%%%%%%%%%%%%%%%%
%%%%%%%%%%%%%%%%%%%%%%%%%%%%%%%%%%%%

\section{Methods}
\textbf{[There's a TON of stuff in here where I'm unsure how much detail to put in the main text vs what to throw into the supplemental material. Eager to hear suggestions!]}
    \begin{itemize}
        \item Introduce MHD equations in $\mathbf{B}$ formulation (unsure how to introduce non-dimensionalization/linearization -- if space weren't a limit I'd first show the dimensional, fully nonlinear equations, then I'd linearize, then non-dimensionalize, but we don't have room for that)
        \item Mention $\mathbf{A}$ formulation, but defer to supplemental materials for full equations
        \item Explain that the whole calculation is separable in $k_x$ and $k_y$ 
        \item Define what we mean by LOP and $G(t_0)$
        \item Broadly summarize that we use the matrix method, cite Squire \& Bhattacharjee, MacTaggart, and I think Reddy
        \item \textbf{Mention sparse vs dense solvers?}
        \item Worth saying something about sanity checks we did in resolution and in number of eigenmodes calculated by sparse solver (if we mention sparse solver)
        \item As explained below in Sec.~[somewhere], in this work we compare two types of calculations, one where we sweep $k_x$ over a range of evenly-spaced values between $0.1$ and $1$ which we refer to as "big box" calculations, and another where we sweep $k_x$ over a range of logarithmically-spaced values between $10^{-3}$ and $1$ which we refer to as "small box" calculations. \textbf{[The following is not actually true right now, but I think I should rerun all my big-box calculations to make it true -- no point in the two different boxes having different kymax]} In both cases, we sweep $k_y$ over a range of evenly-spaced values between $0$ and $1$
        \item Mention that the matrix becomes poorly-conditioned at tiny $k_x$, so while we can't do $k_x = 0$, we're able to do $k_x = 10^{-3}$ and maybe mention the plateau
    \end{itemize}
We use the 3D incompressible MHD equations with finite viscosity and resistivity. We consider small-amplitude disturbances to the equilibrium shear flow $\mathbf{U}_0 = U_0 \tanh(z/d) \hat{\mathbf{e}}_x$ and uniform, streamwise magnetic field $\mathbf{B}_0 = B_0 \hat{\mathbf{e}}_x$. The equations governing these disturbances, linearized about the equilibrium, and non-dimensionalized in terms of the shear layer half-width $d$, the speed of the shear flow $U_0$, and the strength of the magnetic field $B_0$, are then
\begin{equation} \label{eq:lin-momentum}
\frac{\partial}{\partial t} \mathbf{u} + wU' \mathbf{e}_x + U \frac{\partial}{\partial x}\mathbf{u} = - \nabla p + \frac{1}{M_A^2} \frac{\partial}{\partial x}\mathbf{b} + \frac{1}{\mathrm{Re}} \nabla^2 \mathbf{u},
\end{equation}
\begin{equation} \label{eq:lin-induction}
\frac{\partial}{\partial t} \mathbf{b} + U \frac{\partial}{\partial x}\mathbf{b} = \frac{\partial}{\partial x}\mathbf{u} + b_z U' \mathbf{e}_x + \frac{1}{\mathrm{Rm}} \nabla^2 \mathbf{b},
\end{equation}
and
\begin{equation}
\nabla \cdot \mathbf{u} = 0, \qquad \nabla \cdot \mathbf{b} = 0.
\end{equation}
Here, $M_A \equiv U_0/v_A$ is the Alfv\'{e}n Mach number with $v_A \propto B_0$ the Alfv\'{e}n speed, $\mathrm{Re} \equiv U_0 d / \nu$ is the Reynolds number with $\nu$ the kinematic viscosity, and $\mathrm{Rm} \equiv U_0 d / \eta$ is the magnetic Reynolds number with $\eta$ the magnetic diffusivity. Throughout this work, we restrict ourselves to calculations with $\mathrm{Re} = \mathrm{Rm}$, i.e., unity magnetic Prandtl number.
We assume periodic boundary conditions in $x$ and $y$, and no-slip, perfectly conducting boundaries in $z$, as described in the supplemental materials. 

To enforce $\nabla \cdot \mathbf{b} = 0$, we formulate the system in terms of the magnetic vector potential $\mathbf{A}$ in the Coulomb gauge, as described in the supplemental material \textbf{[cite Immie's paper?]}

%We solve the resistive MHD equations linearized about a background field $\Bzero$ and a background shear profile $\Uzero$, cast in terms of magnetic vector potential $\vecpot$ in the Coulomb gauge,
%\begin{equation}
%\rho_0 \left( \frac{\partial}{\partial t} + \mathbf{v} \cdot \nabla \right) \mathbf{v} = -\nabla p + \frac{1}{c} \mathbf{J} \times \mathbf{B} + \rho_0 \nu \nabla^2 \mathbf{v}
%\end{equation}
%\begin{equation}
%\frac{\partial}{\partial t} \mathbf{A} + \nabla \phi = \mathbf{v} \times \mathbf{B} + \eta \nabla^2 \mathbf{A}
%\end{equation}
%\begin{equation}
%\nabla \cdot \mathbf{v} = 0
%\end{equation}
%\begin{equation}
%\nabla \cdot \mathbf{A} = 0.
%\end{equation}
We solve the system using Dedalus \cite{2020PhRvR...2b3068B} and its companion package, eigentools \cite{2021JOSS....6.3079O}.

Gains and linear optimals are calculated as follows. Consider the linear system
\begin{equation}
\frac{\partial U}{\partial t} = \mathcal{L} U,
\label{eq:simplest_system}
\end{equation}
where $U(t)$ is the system state and $\mathcal{L}$ is a time-independent linear differential operator.  
%(Squire \& Bhattacharjee consider the case where $\mathcal{L}$ is time-dependent, but we're not interested in that so I won't consider it here). 
Then the general solution is given by the propagator acting on the initial condition:
\begin{equation}
U(t) = \mathcal{K}(t) U(0),
\label{eq:simplest_solution}
\end{equation}
where
\begin{equation}
\mathcal{K}(t) = \exp(\mathcal{L} t).
\label{eq:simplest_propagator}
\end{equation}
Then, the maximum amplification in terms of some norm $|| \cdot ||_\chi^2$ by time $T$ is
\begin{equation}
G_\chi(T) = \max_{U(0)} \frac{|| \mathcal{K}(T) U(0) ||_\chi^2}{||U(0)||_\chi^2}.
\end{equation}

While several other norms can be considered and might shed insight into the system \textbf{[I know there's a paper from Colm's group to cite here]}, in this work we restrict ourselves to the energy norm, defined as
\begin{equation}
    \text{[Come back to this]}
\end{equation}
[The rest of this paragraph needs to cover Cholesky decomposition, spurious modes, and matrix exponential in eigenmode basis. And it needs to cite Squire \& Bhattacharjee, MacTaggart, and Reddy where appropriate, plus JOSS paper and Boyd's book for spurious modes.] The implementation of this norm is done via a Cholesky decomposition, and the calculation of the matrix exponential is done in the eigenmode basis; both of these methods are described in [cite one/both of the Squire \& Bhattacharjee MRI LOP papers -- also, MacTaggart might cover the latter of those two points too, double-check that, ditto for the Reddy paper]. 
As is standard for systems of this nature \textbf{[Jeff, I don't actually know how broad this phenomenon is, do you?]}, when calculating the eigenmodes of this system, so-called spurious modes often arise [cite Boyd's book]. We separate spurious modes from the rest of the spectrum using the spurious-rejection routine implemented in \cite{2021JOSS....6.3079O} and described in [Boyd's book], and we filter spurious modes from the the calculation of gains and LOPs using the same methods described in [cite MacTaggart, and there's one specific subsection/equation to point to].

Equations \eqref{eq:lin-momentum}-\eqref{eq:lin-induction} are linear, homogeneous in $x$ and $y$, and we use periodic boundary conditions along those axes. Thus, the system can be Fourier transformed in $x$ and $y$, and the resulting equations are entirely separable, i.e., the dynamics at one $(k_x, k_y)$ is entirely decoupled from other wavenumbers. This 

%%%%%%%%%%%%%%%%%%%%%%%%%%%%%%%%%%%%
%%%%%%%%%%%%%%%%%%%%%%%%%%%%%%%%%%%%
%%%%%%%%%%%%%%%%%%%%%%%%%%%%%%%%%%%%

%\section{?Nonmodal growth section?}
\textbf{[Removed the standalone ``Nonmodal growth" section. Things that would have been in here (see commented-out material in the .tex file) should be moved to intro and/or results]}
    % \begin{itemize}
        % \item Possible section giving more detail on different nonmodal mechanisms? I think we want to communicate that uniform magnetic fields allow hydro lift-up to exist without interference (at $k_x = 0$ and without dissipation), and that they introduce their own MHD lift-up. Derivation can go in supplemental material, but we can write down the end results here
        % \item Tie back to stratified KH, since that's what we want to compare our small-box results against
    % \end{itemize}

%%%%%%%%%%%%%%%%%%%%%%%%%%%%%%%%%%%%
%%%%%%%%%%%%%%%%%%%%%%%%%%%%%%%%%%%%
%%%%%%%%%%%%%%%%%%%%%%%%%%%%%%%%%%%%

\section{Small boxes}
    \begin{itemize}
        \item Several aspects of the system are similar to shear layers in neutral fluids, stratified or unstratified:
        \begin{itemize}
            \item In the unstable case, there's a good amount of 3D nonmodal growth ($M_A = 40$ in particular looks a lot like Arratia et al) before modal growth wins at larger $t_0$ and makes things 2D because those are the most-unstable modes
            %\item Both the modal and nonmodal stages at $M_A = 40$ are KE-dominated; the nonmodal stage doesn't look like it's just Orr or lift-up alone
            \item We do get growth even in the stable case! $G \sim 40$ at its peak for $M_A = 1$, $\mathrm{Re} = 250$, for instance
            \item LOPs pretty broadly have some features that look like eddies tilted against the shear , though are typically three-dimensional in nature (this can be tied back to some of the early-90s unstratified shear flow stuff, I think)
            \item In the modally stable cases, waves become more and more excited/present at larger $t_0$'s, and LOPs broaden out at larger $t_0's$
            \item Gain goes down as the system becomes more stable
        \end{itemize}
        \item But a couple things are quite different:
        \begin{itemize}
            \item ME shear production is sometimes preferred over KE shear production -- sometimes KE shear production will be important for a while before ME shear production takes over within the same IVP (see Fig 25 in \texttt{Large\_vs\_small\_boxes/boxes\_notes.tex}), sometimes ME shear production leads the charge from the outset (fig 30 in those notes)
            \item When ME shear production eventually takes over from KE shear production, the phase relationship between $b_x$ and $b_z$ is what you expect from inspecting the energy production terms (see the aside in the MHD Orr scratch notes): while Orr-like KE shear production likes to have $u = -w$, the ME shear production counterpart likes to have $b_x = b_z$ %the t_0=36 IVP for MA=1 is a nice demonstration of this
            \item Except when modal growth is the dominant growth mechanism, the optimal $k_x$ almost always latches onto the minimum allowed, which did not happen in Alexis' work
        \end{itemize}
        \item That last point transitions us into talking about larger boxes...
    \end{itemize}
\section{Big boxes}
    \begin{itemize}
        %\item There was no need to check arbitrarily large boxes in the stratified case because the LOPs weren't constantly latching onto the minimum $k_x$ like they are here
        \item In the stratified case, the LOPs didn't latch on to the minimum $k_x$, but here they do; this motivates us to check larger domains to see how the structure of the LOP varies
        \item We see that gain gets larger and larger as $k_{x, \mathrm{min}}$ gets smaller up to some plateau (although the plateau gets smaller as $t_0$ gets larger), but we find the matrix gets terribly conditioned as $k_x$ gets smaller
        \item We restrict ourselves to $k_x \geq 10^{-3}$ which we find to be a good middle-ground
        \item Even for $M_A \leq 1$ we find HUGE gains in these boxes as long as you're willing to run out to longer $t_0$'s!
        \item In $M_A \leq 1$: strong preference for 2Dyz modes; the LOP is all $v$/$w$ and/or $b_y$/$b_z$, and then the evolved state is all $u$ and/or $b_x$
        \item Pretty clear that hydro lift-up is playing a big role when nonmodal growth is important, often joined by magnetic lift-up when $M_A \leq 1$ -- the importance of lift-up is (I am pretty sure) broadly consistent with some early results for unstratified shear flows without normal-mode instability (like PCF or Poiseuille at subcritical Re), and it's probably worth mentioning that at least in passing
        \item Waves are also a pretty big deal here -- as long as it's modally stable, the LOP seems to swing back and forth between ME vs KE dominated on an Alfvenic timescale (this is something that sets our results apart from the hydro case, though is another connection to the stratified problem; probably can comment on the role of waves in systems that support them)
    \end{itemize}
\section{Discussion}
\begin{itemize}
    \item Although Squire's theorem is known not hold universally in plasma/MHD shear flows [cite Hunt 1966; Hughes \& Tobias], the use of Squire's theorem to motivate 2Dxz studies of KH in plasmas is somewhat routine, and indeed one often finds that 2Dxz modes have the largest growth rates anyway [e.g.~Fraser JFM 2022]. Our results show that, regardless of whether 2Dxz modes have the largest modal growth rates, plasma shear flows in fact often prefer 2Dyz modes, thus 2Dxz simulations likely miss significant growth mechanisms.
    \item Also, astrophysicists are making a huge mistake whenever they say $(M_A > 1) \Leftrightarrow$ (MHD shear flows drive growth)
    \item Furthermore, [big boxes!]
\end{itemize}
Exploiting similarities between stratified shear flows and magnetized ones. Importing idea of linear optimal perturbations to MHD/plasmas. Note that this is beginning of a generalizable idea/method. Can incorporate more ``realistic'' plasma models to see additional physics. Understanding role of subcritical effects, particularly transient growth, on tokamak edge physics is the first step toward understanding optimal control of perturbations [make sure to cite Landreman's work on transient/subcritical growth in tokamak core plasmas]. The ability to control perturbations, and thus shear layer dynamics, could prove pivotal in constructing reliable, efficient fusion power. \textbf{Preceding sentence might be a bit much.}

Talk about how widely astrophysicists use the notion of $M_{A, \mathrm{crit}} \sim 1$ to decide if KH is present, and the importance of KH-induced field amplification in their systems.



\begin{acknowledgments}
We wish to acknowledge people and grants and clusters. Things/people to make sure we acknowledge: the Lux and Alpine clusters; Daniel, Keaton, Bindesh; Pascale's grants as well as Hale/HTMS
\end{acknowledgments}


\bibliography{mhd_linopt}% Produces the bibliography via BibTeX.

\end{document}